\documentclass[a4paper,12pt]{ctexart}
\usepackage{xcolor}
\usepackage{graphicx}
\usepackage[top=1.8cm, bottom=1.8cm, left=1.8cm, right=1.8cm]{geometry}
\usepackage{array}
\usepackage{hyperref}
\usepackage{booktabs}
\hypersetup{colorlinks=true,linkcolor=blue!70!black}
\linespread{1.35}

\title{\textcolor{blue!70!black}{\Huge\textbf{金融学}}}
\subtitle{\textcolor{gray}{经济学门类 · 家长版}}
\author{}
\date{}

\begin{document}
\maketitle
\thispagestyle{empty}

\section{核心数据卡}

\begin{tabular}{ll}
\toprule
专业名称 & 金融学 \\
专业代码 & 020301K \\
所属门类 & 经济学 \\
\midrule
\multicolumn{2}{l}{\textbf{北京选科要求}} \\
\quad 顶尖院校（清北） & 物理必选 \\
\quad 其他院校 & 不限 \\
\midrule
\multicolumn{2}{l}{\textbf{北京代表院校}} \\
\quad 第一梯队 & 北大光华、清华经管、人大财金 \\
\quad 第二梯队 & 央财、对外经贸 \\
\quad 第三梯队 & 首经贸、北工商 \\
\midrule
\multicolumn{2}{l}{\textbf{毕业三年月薪参考（北京）}} \\
\quad 银行总行管培 & 25K--40K \\
\quad 券商投行部 & 40K--80K+ \\
\quad 基金研究员 & 20K--50K \\
\quad 互联网金融 & 25K--45K \\
\bottomrule
\end{tabular}

\section{投入产出比估算}

本科四年学费（公办）约2-3万，生活费和杂费约12-20万。
若读研（2-3年），总投入约30-50万。
毕业后3年累计收入（中位数）约50-80万。
投资回收期约3-5年。

\section{需要留意的风险}

\begin{enumerate}
  \item 学历通胀——十年前本科能进的总行现在要硕士
  \item AI正在替代传统信贷审核和交易岗位
  \item 行业限薪令持续影响薪酬天花板
  \item 竞争极度内卷——清北复交以下无金融的说法虽夸张，但反映行业现实
\end{enumerate}

\section{给家长的建议}

\textbf{成绩一流（清北人财线）}：确保高中数学和英语，目标院校越来越看重物理选科。

\textbf{成绩中上（首经贸/北工商线）}：金融依然是不错的选择，但预期要务实——大概率进银行分支行或保险。建议辅修计算机/法学/外语增加竞争力。

\textbf{数学明显偏弱}：慎重考虑——计量经济学、金融工程、风险管理对数学要求不低。

\end{document}
